\subsection{Formation guidance: NMPC formulation}\label{sec:nmpc}
Each follower is modelled as a point mass in the horizontal plane at constant altitude with state $\mathbf x=[p_x,p_y,V,\chi,\phi,a]^\top$ (north, east, airspeed, course, bank, along-track acceleration) and inputs $\mathbf u=[a_{\mathrm{cmd}},\phi_{\mathrm{cmd}}]^\top$:
\begin{equation}
\dot p_x=V\cos\chi+w_x,\;\dot p_y=V\sin\chi+w_y,\;\dot V=a,\;\dot\chi=\tfrac{g\tan\phi}{V},\;\dot\phi=\tfrac{\phi_{\mathrm{cmd}}-\phi}{\tau_\phi},\;\dot a=\tfrac{a_{\mathrm{cmd}}-a}{\tau_a},
\end{equation}
with $\tau_\phi=\SI{0.25}{s}$, $\tau_a=\SI{0.30}{s}$ and the aircraft's effective wind $\mathbf w_{\mathrm{eff}}$ a first-order lag ($\tau_g=\SI{1}{s}$) of the instantaneous wind, so that a sudden gust accelerates the aircraft over about a second instead of changing its ground velocity instantaneously. The wind consists of a common component (speed $U(2,6)$\,m/s, uniform direction), an individual Ornstein--Uhlenbeck component per aircraft ($\sigma=\SI{0.4}{m/s}$, correlation time \SI{2}{s}) and, in the gust scenarios, a differential lateral doublet applied to the follower only.

\paragraph{Optimal control problem} At \SI{10}{Hz} the NMPC solves, with multiple shooting, RK4 integration, $N=25$ steps of $\Delta t=\SI{0.1}{s}$,
\begin{equation}
\min_{\mathbf u_{0:N-1}}\ \sum_{k=0}^{N-1}\Big(\|\mathbf e_{p,k}\|^2_{Q_p}+\|\mathbf e_{v,k}\|^2_{Q_v}+\|\mathbf u_k\|^2_{R}+\|\mathbf u_k-\mathbf u_{k-1}\|^2_{R_d}\Big)+10\|\mathbf e_{p,N}\|^2_{Q_p}
\end{equation}
subject to the dynamics, $V\in[12,26]$\,m/s, $|\phi|\le35^\circ$, $|a|\le\SI{3}{m/s^2}$, $|a_{\mathrm{cmd}}|\le\SI{3}{m/s^2}$, $|\phi_{\mathrm{cmd}}|\le35^\circ$. Position and velocity references are the delayed, constant-velocity-extrapolated leader state plus the (minimum-jerk) slot trajectory; the wind estimate is a 1.5\,s lag of the true effective wind (an idealised observer). Weights: $Q_p=\mathrm{diag}(1,6)$, $Q_v=\mathrm{diag}(0.4,0.8)$, $R=\mathrm{diag}(0.2,3)$, $R_d=\mathrm{diag}(0.5,8)$ (the lateral channel is weighted more because the wake benefit lies within about $0.1b$). IPOPT \cite{wachter2006} is called through CasADi \cite{andersson2019} with a tolerance of $10^{-4}$, at most 60 iterations and warm starting; the wake term $w_{\mathrm{wake}}\Delta P$ of the preliminary cost is not used. The baseline controller is a PD law in the follower-heading frame ($a_x=0.30e_x+0.9\dot e_x$, $a_y=0.45e_y+1.1\dot e_y$, $\phi_{\mathrm{cmd}}=\arctan(a_y/g)$) with the same limits.

\subsection{High-order barrier-function safety filter}\label{sec:hocbf}
Let $\mathbf r=\mathbf p_f-\mathbf p_L$ and $h(\mathbf x)=\|\mathbf r\|^2-D^2\ge0$ with $D=D_{\mathrm{safe}}+\SI{0.5}{m}$. Differentiating along the follower dynamics, $\dot h=2\mathbf r\cdot\dot{\mathbf r}$ contains no input; $\ddot h=2\|\dot{\mathbf r}\|^2+2\mathbf r\cdot\ddot{\mathbf r}$ contains the acceleration state $a$ and the bank state $\phi$ (through $\dot\chi$) but not $a_{\mathrm{cmd}}$ or $\phi_{\mathrm{cmd}}$; only the third derivative contains the commands, through $\dot a=(a_{\mathrm{cmd}}-a)/\tau_a$ and $\dot\phi=(\phi_{\mathrm{cmd}}-\phi)/\tau_\phi$. The relative degree is therefore $3$ for both inputs (the preliminary version stated 2, which holds only if bank and acceleration are directly commanded). Following \cite{xiao2022} define
\begin{equation}
\psi_0=h,\qquad \psi_i=\dot\psi_{i-1}+\alpha_i\psi_{i-1}\ (i=1,2,3),\qquad (\alpha_1,\alpha_2,\alpha_3)=(1.2,1.6,2.0),
\end{equation}
so that $\psi_3\ge0$ is affine in $\mathbf u$ and renders the set $\{h\ge0\}$ forward invariant \emph{provided} the constraint is feasible within the input bounds. The filter solves at \SI{50}{Hz}
\begin{equation}
\mathbf u^\star=\arg\min_{\mathbf u\in\mathcal U}\tfrac12\|\mathbf u-\mathbf u_{\mathrm{nom}}\|^2\quad\text{s.t.}\quad \mathbf A(\mathbf x)\mathbf u+c(\mathbf x)\ge0,
\end{equation}
where $\mathbf A,c$ are generated symbolically by CasADi from the composition above, using the noisy measured relative position. The QP has two unknowns and one affine constraint over the input box $\mathcal U=[-3,3]\,\mathrm{m/s^2}\times[-35^\circ,35^\circ]$ and is solved by SLSQP in about a millisecond (when the affine barrier half-space conflicts with $\mathcal U$, SLSQP terminates at the box-boundary point minimising constraint violation and clips to $\mathcal U$, rather than introducing an explicit slack variable). Although the filter's internal Lie derivatives include the nominal first-order actuator lags $(\tau_\phi,\tau_a)$ to reach relative degree $3$, the continuous-time forward-invariance guarantee still does not carry over to the closed-loop simulation because of (i) hard input saturation on $\mathcal U$, (ii) unmodelled leader deceleration $\dot{\mathbf v}_L\ne\mathbf 0$ and communication dropouts (the filter differentiates under a constant-velocity leader assumption), (iii) unsteady wind gusts, the effective-wind lag $\tau_g$ and the \SI{1.5}{s} wind-observer lag, (iv) discrete-time \SI{50}{Hz} sampling and the \SI{12}{m/s} airspeed floor, and (v) relative-position measurement noise and GNSS bias. Furthermore, because $h$ is purely radial with fixed class-$\mathcal K$ gains, $\nabla h$ at $(d_x=-3b,d_y=0.85b)$ is predominantly longitudinal, so the myopic single-step QP saturates longitudinal braking before initiating an early lateral break-away manoeuvre of the kind synthesized by backup or adaptive CBFs \cite{chen2021backup,xiao2021ada,molnar2025}; Sec.~\ref{sec:control} therefore reports empirical violation rates rather than a proof.

\subsection{Reconfiguration manoeuvres}\label{sec:reconfdesign}
A \emph{split} moves each follower laterally from its V slot to the centre of its survey zone ($Y=\SI{369.6}{m}$ for the 17-strip plan, \SI{370}{m} for the 14-strip plan; \SI{60}{m} for join-up after take-off and break-up before landing). The lateral reference is a minimum-jerk profile $y(t)=y_0+(y_1-y_0)(10s^3-15s^4+6s^5)$, $s=(t-t_0)/T$, with duration $T=1.875\,Y/v_{\mathrm{lat}}$ for a peak lateral speed $v_{\mathrm{lat}}$ chosen from $\{3,4,6,8,12\}$\,m/s; the NMPC tracks it while keeping the along-track position, so the follower flies a longer path at a higher airspeed (to hold ground speed) and loses the wake benefit as it leaves the slot. The planner selects, for each manoeuvre, the speed that minimises the follower's energy relative to staying in formation, subject to the window using at most half of a one-way transit. Energy over the window is integrated with the full power model using the simulated airspeed, bank and slot offset (the wake term is interpolated from the aerodynamic table), so the cost of a manoeuvre is a simulation output rather than an assumed duration.

\subsection{Mission accounting}\label{sec:accounting}
For each aircraft the mission is a sequence of phases $j$ with power $P_j$ and duration $t_j$; the load-side energy is $E=\sum_jP_jt_j/3600$ and the cell-side energy follows from the battery model. Phases: (i) climb to \SI{120}{m} at \SI{2}{m/s} ($P=P_{\mathrm{cruise}}+mg\,\mathrm{ROC}/\eta_p$); (ii) outbound transit of length $D$ (hypotenuse to the zone for independent aircraft), with the wind-triangle ground speed; (iii) strip legs of \SI{1}{km}; (iv) Dubins turns of length $L_t$ at the turn power; (v) a ferry leg of \SI{1}{km} for an aircraft that ends its last strip on the far side of the polygon; (vi) waiting at the best-endurance speed (\SI{\nnPwait}{W}) until the slowest aircraft finishes, only for strategies that must re-form; (vii) reconfiguration windows; (viii) return transit; (ix) \SI{90}{s} of approach and landing at cruise power. In formation transit the leader role rotates, so each aircraft pays the mean of the leader and follower power, $P_{\mathrm{form}}=(P_{\mathrm{iso}}+2P_{\mathrm{fol}})/3$, with $P_{\mathrm{fol}}$ computed from the drag reduction at the realised slot error (sampled from the control campaign). A rigid V flies the whole 17-strip block with the followers' turn speeds $\Va(1\pm d_y/R_t)$ and no wake benefit in turns. Strip counts per aircraft are $(6,6,5)$ for the 17-strip plan and $(5,5,4)$ for the 14-strip plan.

\subsection{Coverage and acceptance rule}\label{sec:acceptance}
Ground footprints are rectangles $W_g\times H_g$ rotated by the crab angle, triggered every $H_g(1-o_f)$ along track, starting and ending half a footprint outside the polygon. The number of images covering each $\SI{2}{m}$ raster cell is accumulated. A \SI{1}{ha} tile is accepted if at least \SI{99}{\percent} of its cells are covered and at least \SI{95}{\percent} are covered by $\ge5$ images; the preliminary version's additional criteria on checkpoint RMSE \cite{asprs2024} are \emph{not} evaluated here. Strip centres are placed symmetrically; the \SI{58}{\percent} plan is evaluated both with its nominal spacing (\SI{73.92}{m}) and with the smallest spacing that meets the five-ray rule at the block edge (\SI{74.0}{m}).
